\title{DeepSeekMath: Pushing the Limits of Mathematical Reasoning in Open Language Models}

\begin{abstract}
The task of mathematical reasoning remains particularly challenging for language models, primarily due to its inherent complexity and structural requirements. This work presents DeepSeekMath~7B, which is developed by further pre-training DeepSeek-Coder-Base-v1.5 7B using a dataset containing 120B tokens focused on mathematics, curated from Common Crawl, and enriched with both natural language and programming code resources. DeepSeekMath~7B attains a remarkable score of 51.7\% on the rigorous MATH benchmark, performing at a level close to Gemini-Ultra and GPT-4, all without dependency on external toolkits or voting mechanisms. When evaluating self-consistency across 64 sampled outputs, DeepSeekMath~7B reaches 60.9\% on MATH. The enhanced capacity of DeepSeekMath~for mathematical reasoning is primarily attributed to two main aspects: firstly, the extensive exploitation of public web data achieved through a carefully designed data selection process; secondly, the introduction of Group Relative Policy Optimization (GRPO), an adaptation of Proximal Policy Optimization (PPO), which simultaneously bolsters mathematical problem-solving capabilities and improves the efficiency of memory consumption relative to traditional PPO.
\end{abstract}

\section{Introduction}

The emergence of large language models (LLMs) has fundamentally transformed artificial intelligence research in the domain of mathematical reasoning, resulting in notable progress on tasks such as quantitative reasoning benchmarks \citep{MATH} and geometry benchmarks \citep{trinh2024solving}. Furthermore, these systems have demonstrated considerable value in supporting humans tackling challenging mathematical tasks \citep{tao}. Despite these achievements, state-of-the-art models like GPT-4 \citep{gpt4} and Gemini-Ultra \citep{gemini} remain inaccessible to the public, and available open-source alternatives substantially lag behind with respect to performance metrics.

This work presents DeepSeekMath, a specialized language model tailored for mathematics, which demonstrates substantially superior performance over other open-source counterparts and nears GPT-4 in various academic evaluations. Central to this advancement is the construction of the DeepSeekMath Corpus, an extensive and meticulously curated pre-training dataset comprising 120B mathematics-related tokens. The dataset was assembled from Common Crawl (CC) sources via a classifier based on fastText \citep{joulin2016fasttext}. During the initial phase, the classifier was trained using positive samples from OpenWebMath \citep{openwebmath}, supplemented by a wide array of unrelated web content as negative samples. Leveraging this classifier, additional positive mathematical data were identified from CC, then subjected to further curation and validation by human annotators. These newly verified examples were utilized to retrain and refine the classifier for better precision. Evaluation outcomes reveal that the resulting large-scale corpus possesses high quality: DeepSeekMath-Base 7B achieves scores of 64.2\% on GSM8K \citep{gsm8k} and 36.2\% on the advanced MATH benchmark \citep{MATH}, surpassing Minerva 540B \citep{minerva}. Furthermore, the multilingual nature of the DeepSeekMath Corpus translates to observable gains on Chinese mathematical benchmarks \citep{wei2023cmath,agieval}. We view our approach to mathematical data curation as a valuable resource for ongoing research and anticipate further enhancements in subsequent studies.

DeepSeekMath-Base is initialized from DeepSeek-Coder-Base-v1.5 7B \citep{deepseek-coder}, since we found that leveraging a code-focused pretrained model yields superior results compared to initializing from a standard large language model. Notably, the math-centric training further boosts performance not only in mathematics but also enhances general reasoning, as demonstrated by gains on the MMLU \citep{mmlu} and BBH \citep{bbh} benchmarks.

Following this pre-training stage, we perform mathematical instruction tuning on DeepSeekMath-Base utilizing datasets for chain-of-thought \citep{cot}, program-of-thought \citep{pot,pal}, and tool-integrated reasoning \citep{tora}. The resultant DeepSeekMath-Instruct 7B model surpasses all other models in the 7B parameter class and attains performance on par with state-of-the-art open-source instruction-tuned models with 70B parameters.

Additionally, we present Group Relative Policy Optimization (GRPO), a variant of the Proximal Policy Optimization (PPO) algorithm \citep{schulman2017proximal} tailored for reinforcement learning (RL). GRPO eliminates the need for a critic model and instead derives the baseline using group-wise scores, thereby offering notable computational efficiency. Using only a portion of the English instruction tuning corpus, GRPO delivers significant enhancements over the already robust DeepSeekMath-Instruct model in both domain-specific (GSM8K: 82.9\% $\rightarrow$ 88.2\%, MATH: 46.8\% $\rightarrow$ 51.7\%) and out-of-domain mathematics tasks (e.g., CMATH: 84.6\% $\rightarrow$ 88.8\%) throughout the RL stage. Moreover, we establish a unified framework that interprets a range of techniques—including Rejection Sampling Fine-Tuning (RFT) \citep{yuan2023scaling}, Direct Preference Optimization (DPO) \citep{dpo}, PPO, and GRPO—as variants of direct or simplified RL strategies. Through comprehensive experimentation—such as contrasting online and offline training, outcome versus process-based supervision, and single-turn versus iterative RL—we rigorously analyze the crucial components of this unified framework. Finally, we elucidate the mechanisms by which our RL approach enhances the performance of instruction-tuned models and outline prospective pathways for achieving even more effective RL within this coherent paradigm.

\subsection{Contributions}
This work presents advancements in large-scale mathematical pre-training, in addition to a thorough investigation into the application and assessment of reinforcement learning.

\textbf{Math Pre-Training at Scale}

\begin{itemize}[topsep=0pt]
    \item We demonstrate that publicly available Common Crawl data is a rich resource for mathematics-related information. Leveraging a carefully crafted data selection pipeline, we have assembled the DeepSeekMath Corpus: a high-quality mathematical dataset containing 120B tokens sourced from web pages filtered for math relevance. This corpus is nearly 7 times larger than the dataset of math web pages used by Minerva \citep{minerva} and is approximately 9 times bigger than the newly published OpenWebMath \citep{openwebmath}.

    \item Our DeepSeekMath-Base 7B pre-trained model attains a level of mathematical reasoning performance on par with Minerva 540B \citep{minerva}, suggesting that model size is not the sole determinant of success in mathematical tasks. High-quality data and effective pre-training can enable smaller models to exhibit strong mathematical reasoning capabilities.

    \item We report observations from our mathematical training experiments, noting that pre-training on code before engaging in math-focused training enhances model proficiency in mathematical problem solving, regardless of tool assistance. This finding contributes to the debate over whether code pre-training improves general reasoning, providing evidence that it benefits at least mathematical reasoning.

    \item Despite the prevalence of training on arXiv papers within math-focused research, our experiments reveal that incorporating arXiv data yields no significant gains on any of the mathematical benchmarks considered in this study.
\end{itemize}

\textbf{Exploration and Analysis of Reinforcement Learning}

\begin{itemize}[topsep=0pt]
    \item We present Group Relative Policy Optimization (GRPO), a novel reinforcement learning algorithm that prioritizes both efficiency and effectiveness. Unlike traditional methods, GRPO dispenses with the critic model and instead leverages group scores to approximate the baseline, thereby substantially cutting down on the computational requirements typically needed for Proximal Policy Optimization (PPO).
    \item We show that GRPO markedly boosts the capabilities of our instruction-tuned model, DeepSeekMath-Instruct, using only instruction-tuning datasets. Additionally, we find that reinforcement learning with GRPO leads to notable improvements in the model’s generalization to tasks outside the original domain.
    \item We introduce a unified framework that contextualizes diverse approaches such as RFT, DPO, PPO, and GRPO. Comprehensive experiments—including analyses contrasting online and offline learning, outcome versus process supervision, and single-step versus iterative reinforcement learning—are conducted to systematically elucidate the core aspects underpinning this unified perspective.
    \item Leveraging this unified view, we delve into the foundational drivers behind reinforcement learning's success and outline several promising directions for further optimizing the reinforcement learning of large language models (LLMs).
\end{itemize}

\subsection{Summary of Evaluations and Metrics}

\begin{itemize}[topsep=0pt]
    \item \textbf{English and Chinese Mathematical Reasoning}: We rigorously evaluate our models using a diverse set of benchmarks in both English and Chinese, spanning mathematical tasks from elementary through collegiate difficulty. English benchmarks comprise GSM8K \citep{gsm8k}, MATH \citep{MATH}, SAT \citep{llemma}, OCW Courses \citep{minerva}, and MMLU-STEM \citep{mmlu}; Chinese datasets encompass MGSM-zh \citep{mgsm}, CMATH \citep{wei2023cmath}, Gaokao-MathCloze \citep{agieval}, and Gaokao-MathQA \citep{agieval}. The assessment focuses on the models’ ability to independently produce complete text-based solutions and to solve tasks via Python code.
    
    On the English datasets, DeepSeekMath-Base demonstrates competitiveness with the proprietary Minerva 540B \citep{minerva}, while exceeding all other open-source models (such as Mistral 7B \citep{mistral} and Llemma-34B \citep{llemma}) by considerable margins, irrespective of previous mathematical pre-training. Remarkably, DeepSeekMath-Base also achieves leading results on Chinese datasets, potentially attributable to its incorporation of high-quality multilingual math pre-training data, diverging from previous work \citep{minerva,llemma} that utilized only English sources. Upon subsequent instruction tuning and reinforcement learning, DeepSeekMath-Instruct and DeepSeekMath-RL exhibit outstanding performance—surpassing the 50\% accuracy threshold on the challenging, competition-level MATH dataset for the first time among open-source efforts.
\end{itemize}

\item \textbf{Formal Mathematics}: DeepSeekMath-Base is assessed using the informal-to-formal theorem proving task described in \citep{dsp_proof}, employing the miniF2F benchmark \citep{minif2f} and selecting Isabelle \citep{isabelle} as the proof assistant. The model exhibits robust few-shot performance in the autoformalization setting.

\item \textbf{Natural Language Understanding, Reasoning, and Code}: To thoroughly characterize the models’ abilities in comprehension, reasoning, and programming, we benchmark DeepSeekMath-Base on Massive Multitask Language Understanding (MMLU) \citep{mmlu}, a suite comprising 57 diverse multiple-choice tasks; BIG-Bench Hard (BBH) \citep{bbh}, featuring 23 demanding tasks primarily centered on multi-step reasoning; as well as on HumanEval \citep{codex} and MBPP \citep{mbpp}, both established for code evaluation. Pre-training with mathematical data leads to improvements in both language comprehension and reasoning capabilities.

\section{Math Pre-Training}

\subsection{Data Collection and Decontamination} 

This section details the methodology for building the DeepSeekMath Corpus from Common Crawl. As shown in Figure \ref{fig:data_collect}, we describe an iterative pipeline for efficiently compiling a large-scale mathematical dataset from Common Crawl, starting with a seed collection—namely, a small, high-quality body of math-relevant data. Notably, this pipeline can be adapted for other specialized domains such as programming.

Initially, we select OpenWebMath \citep{openwebmath}—a well-curated dataset of mathematical web resources—as our starting point. We leverage this dataset to train a fastText model \citep{joulin2016fasttext} that retrieves additional web pages with similar mathematical content. For training, 500,000 samples from the seed corpus are randomly sampled as positive instances, and another 500,000 web pages from Common Crawl serve as negatives. We employ the open-source fastText library\footnote{\url{https://fasttext.cc}}, configuring a vector size of 256, a learning rate of 0.1, a maximum n-gram length of 3, a minimum word occurrence threshold of 3, and three training epochs. The Common Crawl dataset is initially reduced in size using both URL-based deduplication and near-deduplication, yielding a corpus of 40B HTML documents. Using the fastText classifier, we extract math-related web pages from this deduplicated subset. To eliminate low-quality mathematical documents, pages are sorted according to their fastText scores, with only the highest-scoring examples retained. Pre-training experiments are then conducted using the top 40B, 80B, 120B, and 160B tokens to determine the optimal data volume. In the initial pass, we retain only the top 40B tokens.

Following the initial round of data collection, a significant portion of mathematical web pages remains uncollected. This gap largely stems from the fastText model’s limited exposure to sufficiently varied positive samples. To remedy this, we expand the seed corpus by incorporating additional mathematical web sources, thereby aiming to enhance the fastText model’s performance. Our strategy begins by dividing the entirety of the Common Crawl dataset into separate domains, where each domain encompasses web pages under a common base URL. For every domain, we compute the fraction of web pages that were successfully retrieved during the first iteration. Domains in which more than 10\% of pages were collected are categorized as being math-centric (for example, \url{mathoverflow.net}).

We then proceed to manually curate and annotate URLs within these math-centric domains that are directly linked to mathematical content (such as \url{mathoverflow.net/questions}). Uncollected web pages under these curated URLs are subsequently appended to the seed set, enabling us to accumulate a more comprehensive set of positive training examples. This process leads to the training of a more effective fastText model, capable of retrieving a wider range of mathematical data in subsequent collection rounds. Through four such iterations, we ultimately acquire a corpus of 35.5M mathematical web pages, amounting to 120B tokens. Notably, by the fourth round, approximately 98\% of the data had already been gathered in the prior iteration, prompting us to discontinue further collection.

To ensure that our corpus remains free from test data contamination, we employ the filtering methodology outlined in \cite{deepseek-coder}: specifically, web pages are excluded if they contain questions or answers drawn from established mathematical benchmarks in English (e.g., GSM8K \citep{gsm8k}, MATH \citep{MATH}) or Chinese (e.g., CMATH \citep{wei2023cmath}, AGIEval \citep{agieval}). Our filtering protocol operates as follows: any web page is eliminated if any segment of text contains a 10-gram phrase exactly matching a substring present in any of the aforementioned benchmarks. For benchmark passages shorter than 10 grams but containing at least 3 grams, we conduct exact string matching to remove any potentially contaminated pages from our math corpus.

\subsection{Validating the Quality of the DeepSeekMath~Corpus}

We conduct pre-training evaluations to assess the performance of the DeepSeekMath~Corpus in relation to several recently released corpora designed for mathematical training:

\begin{itemize}[topsep=0pt]
    \item \textbf{MathPile} \citep{mathpile}: a composite corpus containing 8.9B tokens, sourced from various repositories such as textbooks, Wikipedia, ProofWiki, CommonCrawl, StackExchange, and predominantly arXiv (over 85\% of the corpus);
    \item \textbf{OpenWebMath} \citep{openwebmath}: a collection of CommonCrawl material specifically filtered for mathematical content, comprising 13.6B tokens;
    \item \textbf{Proof-Pile-2} \citep{llemma}: a mathematical dataset integrating OpenWebMath, AlgebraicStack (offering 10.3B tokens of mathematical code), and arXiv documents (28.0B tokens). For experiments on Proof-Pile-2, we adhere to the 2:4:1 arXiv:Web:Code ratio specified in \cite{llemma}.
\end{itemize}

\subsubsection{Training Setting}
\label{sec:quality-policy}
We perform math-specific training on a general-purpose language model containing 1.3B parameters, which utilizes the same architecture as DeepSeek LLMs \citep{deepseek-llm}, henceforth referred to as DeepSeek-LLM 1.3B. Separate models are trained on each distinct mathematical dataset for a total of 150B tokens. All training is carried out using the streamlined and resource-efficient HAI-LLM \citep{haillm} infrastructure. Adhering to the optimization procedures established by DeepSeek LLMs, the AdamW optimizer \citep{adamW} is employed with $\beta_1=0.9$, $\beta_2=0.95$, and a $\mathrm{weight\_decay}$ of 0.1. The learning rate follows a multi-stage schedule: it linearly increases during the first 2,000 warmup steps to reach its maximum, then decays to 31.6\% of this maximum after 80\% of total training, and is further reduced to 10.0\% of the peak after 90\% of training is completed. The highest learning rate is set to 5.3e-4, the batch size is configured to process 4M tokens, and the context window is set at 4K tokens.

\subsubsection{Evaluation Results}

\textbf{The DeepSeekMath~Corpus excels in data quality, encompasses mathematical data across multiple languages, and surpasses other corpora in scale.}

\begin{itemize}[topsep=0pt]
    \item \textbf{High-quality}:
    To assess performance, we utilize 8 different mathematical benchmarks under few-shot chain-of-thought prompting \cite{cot}. As indicated in Table \ref{tab:corpora_comparison}, models trained on the DeepSeekMath~Corpus consistently outperform others. Further, Figure \ref{fig:corpora_comparisons} highlights that, even at 50B tokens (which equates to a complete pass through Proof-Pile-2), the DeepSeekMath-trained model surpasses its counterpart, underscoring superior average data quality within the DeepSeekMath~Corpus.
    \item \textbf{Multilingual}:
    Data within the DeepSeekMath~Corpus spans several languages, with English and Chinese being most prevalent. Table \ref{tab:corpora_comparison} reveals that using DeepSeekMath~Corpus leads to gains in mathematical reasoning capabilities in both English and Chinese. Conversely, alternative corpora, which are largely restricted to English, fail to offer such benefits and may negatively impact performance on Chinese tasks.
    \item \textbf{Large-scale}:
    The size of the DeepSeekMath~Corpus is several multiples greater than that of previous mathematical datasets. As shown in Figure \ref{fig:corpora_comparisons}, DeepSeek-LLM 1.3B achieves faster and more sustained performance growth when trained on DeepSeekMath~Corpus. By comparison, other corpora, being much smaller and reused several times during training, yield models whose accuracy quickly stagnates.
\end{itemize}

\subsection{Training and Evaluating DeepSeekMath-Base 7B}

Here, we present DeepSeekMath-Base 7B, a foundational model notable for its advanced mathematical reasoning skills. This model is based on DeepSeek-Coder-Base-v1.5 7B \citep{deepseek-coder} and is trained on 500B tokens. The composition of the training dataset is as follows: 56\% DeepSeekMath Corpus, 4\% AlgebraicStack, 10\% arXiv, 20\% Github code, and 10\% Common Crawl natural language data in both English and Chinese. The training configuration closely follows the approach in Section \ref{sec:quality-policy}, with the exception that the learning rate's peak is adjusted to 4.2e-4 and the batch size is set to 10M tokens.

We systematically evaluate DeepSeekMath-Base 7B on a range of tasks to measure its mathematical competence. These include generating complete mathematical solutions independently of external resources, leveraging tools to address math problems, and executing formal theorem verification. Additionally, we broaden our evaluation to encompass natural language understanding, logical reasoning, and programming proficiency, thereby establishing a comprehensive overview of the model’s baseline capabilities.

\paragraph{Mathematical Problem Solving with Step-by-Step Reasoning}

To gauge DeepSeekMath-Base's problem-solving aptitude, we utilize few-shot chain-of-thought prompting \citep{cot} across eight distinct benchmarks in both English and Chinese. The chosen benchmarks cover quantitative reasoning tasks (such as GSM8K \citep{gsm8k}, MATH \citep{MATH}, and CMATH \citep{wei2023cmath}) as well as multiple-choice assessments (including MMLU-STEM \citep{mmlu} and Gaokao-MathQA \citep{agieval}). These benchmarks collectively represent a wide spectrum of mathematical domains, ranging from elementary concepts to problems at the college level.

As detailed in Table \ref{tab:base_math_cot}, DeepSeekMath-Base 7B achieves the highest scores among open-source base models across all eight evaluated benchmarks. This group includes not only the highly popular Mistral 7B \citep{mistral} but also Llemma 34B \citep{llemma}, a recent release that benefits from math-focused pretraining on Proof-Pile-2 \citep{llemma}. On the challenging, competition-level MATH dataset, DeepSeekMath-Base establishes a new standard, exceeding the best open-source base model results by more than 10 percentage points, and even surpasses Minerva 540B \citep{minerva}—a proprietary model built upon PaLM \citep{palm}, trained on mathematical corpora, and possessing over 77 times as many parameters.

\paragraph{Mathematical Problem Solving with Tool Use} 

We assess models' capabilities in program-based mathematical reasoning on the GSM8K and MATH benchmarks, leveraging few-shot program-of-thought prompting methods \citep{pot,pal}. The models are instructed to approach problems by generating Python scripts, utilizing libraries such as \textit{math} and \textit{sympy} to handle sophisticated calculations, with the final program output being taken as the model’s solution. Table \ref{tab:base_math_pot_proof} reveals that DeepSeekMath-Base 7B sets a new performance bar, outperforming the previous leading model, Llemma 34B.

\paragraph{Formal Mathematics}

The automation of formal mathematical proofs is increasingly valued for ensuring correctness, reliability, and improved productivity in mathematical practice. We test DeepSeekMath-Base 7B using the informal-to-formal proof generation benchmark from \citep{dsp_proof}, which requires the transformation of an informal statement (with its formal translation and an informal proof) into a valid formal proof. The evaluation utilizes miniF2F \citep{minif2f}, a formal benchmark reflecting Olympiad-level mathematical problem solving; formal proofs are constructed in Isabelle using few-shot prompting. In line with the methodology of \cite{dsp_proof}, our approach has models output proof sketches, then employs the Sledgehammer automated prover \citep{sledgehammer} to complete any unresolved proof steps. Results in Table \ref{tab:base_math_pot_proof} underscore DeepSeekMath-Base 7B’s competitive proficiency in proof formalization automation.

\paragraph{Natural Language Understanding, Reasoning, and Code}

For broader evaluation, we measure model performance in natural language understanding (MMLU \citep{mmlu}), complex reasoning (BBH \citep{bbh}), and programming (HumanEval \citep{codex}; MBPP \citep{mbpp}). Table \ref{tab:understanding_reasoning_code} shows that DeepSeekMath-Base 7B not only achieves marked gains over its predecessor, DeepSeek-Coder-Base-v1.5 \citep{deepseek-coder}, on MMLU and BBH, but also demonstrates that focused mathematical training enhances language comprehension and reasoning. Furthermore, by continuing training with code-related tokens, DeepSeekMath-Base 7B successfully preserves the strong coding abilities of DeepSeek-Coder-Base-v1.5, maintaining high scores on both code benchmarks. Collectively, these results indicate that DeepSeekMath-Base 7B robustly surpasses the general-purpose Mistral 7B model \citep{mistral} across all three categories: reasoning, coding, and understanding.

\section{Supervised Fine-Tuning}

\subsection{SFT Data Curation}

We assembled a dataset for mathematical instruction tuning that spans both English and Chinese problems, representing a broad spectrum of mathematical disciplines and complexity levels. Each problem is matched with detailed solutions formatted in chain-of-thought (CoT) \citep{cot}, program-of-thought (PoT) \citep{pot,pal}, or tool-integrated reasoning formats \citep{tora}. The dataset comprises a total of 776,000 training examples.

\begin{itemize}[topsep=0pt]
    \item \textbf{English mathematical datasets}: We enrich the GSM8K and MATH collections by providing tool-integrated solution annotations, and further incorporate a selected subset of MathInstruct \citep{MathInstruct} and the Lila-OOD training dataset \citep{lila}, in which solutions leverage CoT or PoT methodologies. The English segment addresses a variety of mathematical areas, such as algebra, probability, number theory, calculus, and geometry.
    \item \textbf{Chinese mathematical datasets}: Our compilation includes Chinese K-12 mathematics problems distributed across 76 sub-domains—for example, linear equations—with annotations supplied in both CoT and tool-integrated formats.
\end{itemize}

\subsection{Training and Evaluating DeepSeekMath-Instruct 7B}

Here, we present DeepSeekMath-Instruct 7B, a model refined through mathematical instruction fine-tuning that builds upon DeepSeekMath-Base. During training, individual samples are concatenated at random to fill a sequence up to 4K tokens in length. The optimization process involves 500 steps, utilizing a batch size of 256 and maintaining a fixed learning rate of 5e-5.

To assess the mathematical reasoning ability of our models—with and without access to external tools—we evaluate on four quantitative reasoning benchmarks across English and Chinese languages. For comparative analysis, we include state-of-the-art models available at the time:

\begin{itemize}[topsep=0pt]
    \item \textbf{Closed-source models}: The baseline includes: (1) GPT family models, with GPT-4 \citep{gpt4} and GPT-4 Code Interpreter~\footnote{\url{https://openai.com/blog/chatgpt-plugins\#\#code-interpreter}} as top-performing examples; (2) Gemini Ultra and Pro \citep{gemini}; (3) Inflection-2 \citep{inflection-2}; (4) Grok-1~\footnote{\url{https://x.ai/model-card}}; as well as recent entries by Chinese technology firms: (5) Baichuan-3~\footnote{\url{https://www.baichuan-ai.com}}; and (6) the latest member of the GLM family, GLM-4~\footnote{\url{https://open.bigmodel.cn/dev/api\#glm-4}} \citep{glm}.
\end{itemize}

These models serve general-purpose applications and have typically been subjected to various alignment processes.
    \item \textbf{Open-source models} encompass: broad models such as (1) DeepSeek-LLM-Chat 67B \citep{deepseek-llm}, (2) Qwen 72B \citep{qwen}, (3) SeaLLM-v2 7B \citep{seallm}, and (4) ChatGLM3 6B \citep{chatglm3}. This category also includes models with explicit mathematical enhancements: (5) InternLM2-Math 20B~\footnote{\url{https://github.com/InternLM/InternLM-Math}}, derived from InternLM2 with subsequent mathematics-specific pretraining and instruction tuning, (6) Math-Shepherd-Mistral 7B, which employs PPO training \citep{schulman2017proximal} on Mistral 7B \citep{mistral} utilizing a process-supervised reward framework, (7) the WizardMath series \citep{wizardmath}, which elevates mathematical reasoning capabilities for Mistral 7B and Llama-2 70B \citep{llama2} via evolve-instruct—a form of instruction tuning that leverages AI-generated instructions—and PPO training, mainly leveraging problems from GSM8K and MATH, (8) MetaMath 70B \citep{metamath}, which fine-tunes Llama-2 70B with an extended dataset including GSM8K and MATH, (9) ToRA 34B \cite{tora}, a CodeLlama 34B variant adapted for tool-augmented mathematical reasoning, and (10) MAmmoTH 70B \citep{MathInstruct}, where Llama-2 70B undergoes instruction tuning on MathInstruct.

As detailed in Table \ref{tab:sft_rl_math}, in evaluations prohibiting tool-assisted reasoning, DeepSeekMath-Instruct 7B exhibits exceptional step-by-step problem-solving performance. Specifically, on the MATH competition dataset, our model achieves at least a 9\% absolute improvement over all open-source models as well as most closed-source systems (e.g., Inflection-2, Gemini Pro), including those with significantly larger model sizes (e.g., Qwen 72B) or with mathematics-oriented reinforcement learning enhancements (e.g., WizardMath-v1.1 7B). Although DeepSeekMath-Instruct attains parity with the proprietary Chinese models GLM-4 and Baichuan-3 on MATH, it falls short of the performance of GPT-4 and Gemini Ultra.

In scenarios where models are permitted to utilize both natural language reasoning and programmatic tool use for solving problems, DeepSeekMath-Instruct 7B attains nearly 60\% accuracy on the MATH dataset, exceeding the accuracy of all previously published open-source models. Across other evaluation datasets, it remains competitive with DeepSeek-LLM-Chat 67B, the prior state-of-the-art, which is an order of magnitude larger.

\section{Reinforcement Learning}

\subsection{Group Relative Policy Optimization} The application of reinforcement learning (RL) has demonstrated marked improvements in LLMs’ mathematical reasoning following the Supervised Fine-Tuning (SFT) phase \citep{wang2023math,wizardmath}. Here, we present Group Relative Policy Optimization (GRPO), our proposed efficient RL method.

\subsubsection{From PPO to GRPO} Proximal Policy Optimization (PPO) \citep{schulman2017proximal} stands as a foundational actor-critic reinforcement learning approach, widely adopted during the RL-based fine-tuning of LLMs \citep{ouyang2022training}. It functions by maximizing the following surrogate objective function:

\begin{equation}
\footnotesize
    \mathcal{J}_{PPO}(\theta) = \mathbb{E}{[q \sim P(Q), o \sim \pi_{\theta_{old}}(O|q)]} \frac{1}{|o|} \sum_{t=1}^{|o|} \min \left[ \frac{\pi_\theta(o_{t} | q, o_{<t})}{\pi_{\theta_{old}}(o_{t} | q, o_{<t})} A_{t}, \text{clip} \left( \frac{\pi_\theta(o_{t} | q, o_{<t})}{\pi_{\theta_{old}}(o_{t} | q, o_{<t})}, 1 - \varepsilon, 1 + \varepsilon \right)  A_{t} \right] ,
\end

In this context, $\pi_{\theta}$ and $\pi_{\theta_{old}}$ denote the present and preceding policy models, while $q$ and $o$ represent the questions and the sampled outputs, respectively—drawn from both the question dataset and from the earlier policy $\pi_{\theta_{old}}$. The parameter $\varepsilon$ serves as a hyper-parameter for clipping, as introduced in PPO to ensure stable training. The term $A_t$ indicates the advantage, derived by employing Generalized Advantage Estimation (GAE) \citep{gae}, which utilizes rewards $\{r_{\ge t}\}$ along with a learned value function $V_{\psi}$. Consequently, PPO requires simultaneous training of a value function in addition to the policy model. To prevent excessive optimization towards the reward model, the widely adopted technique is to incorporate a token-wise KL penalty between the current model and a reference model at every token position within the reward signal \citep{ouyang2022training}, specified as

\begin{equation}
    r_{t} = r_\varphi(q, o_{\le t}) - \beta \log\frac{\pi_{\theta}(o_{t}|q, o_{<t})}{\pi_{ref}(o_{t}|q, o_{<t})},
\label{eq:PPO-reward}
\end{equation}

where $r_\varphi$ stands for the reward model, $\pi_{ref}$ corresponds to the reference model—often instantiated as the original SFT model—and $\beta$ indicates the scaling factor for the KL penalty.

Due to the fact that the value function in PPO is usually a separate neural network of similar size to the policy model, this approach imposes considerable demands on both memory and computation. Furthermore, during reinforcement learning, the value function acts as a baseline to reduce the variance of the advantage estimates. However, within LLM applications, reward models are generally used to score only the final token in the sequence, making it more difficult to train an effective value function that accurately assesses every individual token. To overcome these challenges, as visualized in Figure \ref{fig:grpo}, we introduce Group Relative Policy Optimization (GRPO). GRPO removes the necessity for explicit value function estimation as required in PPO; instead, it takes the mean reward of several outputs—sampled in response to the same prompt—as a baseline. In particular, for a question $q$, GRPO samples a batch of responses $\{o_1, o_2, \cdots, o_G\}$ from the previous policy $\pi_{\theta_{old}}$ and updates the policy model by maximizing the following objective:

\begin{equation}
\footnotesize
\begin{split}
    \mathcal{J}_{GRPO}(\theta) &= \mathbb{E}{[q \sim P(Q), \{o_i\}_{i=1}^G \sim \pi_{\theta_{old}}(O|q)]}  \\
    & \frac{1}{G}\sum_{i=1}^G\frac{1}{|o_i|} \sum_{t=1}^{|o_i|} \left\{ \min \left[ \frac{\pi_\theta(o_{i,t} | q, o_{i,<t})}{\pi_{\theta_{old}}(o_{i,t} | q, o_{i,<t})} \hat{A}_{i,t}, \text{clip} \left( \frac{\pi_\theta(o_{i,t} | q, o_{i,<t})}{\pi_{\theta_{old}}(o_{i,t} | q, o_{i,<t})}, 1 - \varepsilon, 1 + \varepsilon \right)  \hat{A}_{i,t} \right] - \beta \mathbb{D}_{KL}\left[\pi_{\theta} || \pi_{ref}\right]\right\} ,
\end{split}
\label{eq:GRPO-obj}
\end{equation}

Here, both $\varepsilon$ and $\beta$ are tunable hyper-parameters. The term $\hat{A}_{i,t}$ represents the advantage, but is specifically computed with reference to the relative rewards among the outputs within each group. Further details on this computation are given in the following subsections. The group-relative computation of advantages employed by GRPO harmonizes naturally with the pairwise comparative framework often used to train reward models, since such models are generally developed from datasets consisting of output comparisons for identical prompts. Another key distinction is that, rather than integrating the KL penalty into the reward as in PPO, GRPO imposes KL regularization by directly including the divergence between the updated policy and the reference policy within the objective function

\begin{algorithm}[t]
  \small
  \caption{Iterative Group Relative Policy Optimization}
  \textbf{Input} initial policy model $\pi_{\theta_{\text{init}}}$; reward models $r_\varphi$; task prompts $\mathcal{D}$; 
  hyperparameters $\varepsilon$, $\beta$, $\mu$
  \begin{algorithmic}[1]
    \State Assign the initial policy model: $\pi_\theta \leftarrow \pi_{\theta_{\text{init}}}$
    \For{iteration = 1, \dots, I}
       \State Set the reference model: $\pi_{ref} \leftarrow \pi_{\theta}$
      \For{step = 1, \dots, M}
      \State Draw a batch $\mathcal{D}_b$ from $\mathcal{D}$
      \State Synchronize the old policy model: $\pi_{\theta_{old}} \leftarrow \pi_{\theta}$ 
      \State For every question $q \in \mathcal{D}_b$, generate $G$ outputs $\{o_i\}_{i=1}^G \sim \pi_{\theta_{old}} (\cdot \mid q) $
      \State Evaluate each $o_i$ using $r_{\varphi}$ to obtain rewards $\{r_i\}_{i=1}^{G}$
      \State Compute token-level group relative advantages $\hat{A}_{i,t}$ for each $t$-th token in $o_i$
      \For{GRPO iteration = 1, \dots, $\mu$}
        \State Improve the policy $\pi_{\theta}$ by maximizing the GRPO objective (Equation \ref{eq:GC-GRPO})
      \EndFor
    \EndFor \State Continuously update $r_\varphi$ by replay-based training. 
    \EndFor 
  \end{algorithmic}
  \textbf{Output} $\pi_\theta$
  \label{alg:iter-grpo}
\end{algorithm}

\subsubsection{Outcome Supervision RL with GRPO} More precisely, for each question $q$, a collection of candidate outputs $\{o_1, o_2, \cdots, o_G\}$ is drawn from the current old policy model $\pi_{\theta_{old}}$. These outputs are each assigned a scalar reward using a reward model, giving a set $\mathbf{r}=\{r_1, r_2, \cdots, r_G\}$. The obtained rewards are then standardized by subtracting their group mean and dividing by their standard deviation. In the setting of outcome supervision, the normalized reward at the conclusion of every output $o_i$ is used as the advantage $\hat{A}_{i, t}$ for all positions $t$ within $o_i$, namely, $\hat{A}_{i, t} = \widetilde{r}_i = \frac{r_i- {\rm mean}(\mathbf{r})}{{\rm std}(\mathbf{r})}$. The policy is then trained to maximize the criterion described in equation (\ref{eq:GRPO-obj}).

\subsubsection{Process Supervision RL with GRPO} Relying solely on end-of-trajectory outcome rewards may not provide fine-grained or sufficient feedback for policies tackling complex mathematical reasoning tasks. In line with the methodology of \cite{wang2023math}, this work also considers process supervision, in which rewards are delivered at each intermediate reasoning step. Formally, after generating $G$ sampled outputs $\{o_1, o_2, \cdots, o_G\}$ for question $q$, a process reward model assigns rewards for each step within each output, yielding $\mathbf{R} = \{ \{r_1^{index(1)},\cdots,r_1^{index(K_1)}\}, \cdots,  \{r_G^{index(1)},\cdots,r_G^{index(K_G)}\} \}$, where $index(j)$ identifies the ending token of the $j$-th reasoning step and $K_i$ denotes the number of reasoning steps in $o_i$. These stepwise rewards are normalized as $\widetilde{r}_i^{index(j)} = \frac{r_i^{index(j)} - {\rm mean}(\mathbf{R})}{{\rm std}(\mathbf{R})}}$. For each token $t$, the corresponding advantage is computed as the sum over normalized rewards from the current and all subsequent steps: $\hat{A}_{i, t} = \sum_{index(j) \ge t} \widetilde{r}_i^{index(j)}$. Training then seeks

\subsubsection{Iterative RL with GRPO} As the reinforcement learning process advances, the previously trained reward model may no longer provide adequate supervision for the updated policy model. To address this, we investigate an iterative RL approach using GRPO. As detailed in Algorithm \ref{alg:iter-grpo}, iterative GRPO involves producing fresh reward model training datasets derived from samples generated by the current policy model. The reward model is then continually refined using a replay mechanism, retaining 10\% of historical data in the training buffer. Subsequently, the policy model is set as the new reference model and is further optimized with supervision from the latest version of the reward model.

\subsection{Training and Evaluating DeepSeekMath-RL}

Our reinforcement learning experiments utilize DeepSeekMath-Instruct 7B as the backbone. The RL training set consists of approximately 144K chain-of-thought questions relevant to GSM8K and MATH, selected from the SFT data. We intentionally omit other types of SFT questions during RL to assess how RL training influences benchmarks without corresponding in-distribution data. Following the procedure outlined by \citep{wang2023math}, we assemble the training data for the reward model. The reward model initialization is performed with DeepSeekMath-Base 7B using a 2e-5 learning rate. For the GRPO algorithm, the policy model's learning rate is configured to 1e-6, and the KL penalty coefficient is set at 0.04. During training, $64$ outputs are sampled for each question. The maximum sequence length is capped at 1024 tokens, with a batch size of 1024. After every exploration phase, the policy model undergoes a single update. Evaluation of DeepSeekMath-RL 7B is performed on the same benchmarks as DeepSeekMath-Instruct 7B. For DeepSeekMath-RL 7B, tasks associated with GSM8K and MATH in chain-of-thought format are considered in-domain, while remaining benchmarks are classified as out-of-domain.

Table \ref{tab:sft_rl_math} presents a comparative analysis of open-source and closed-source models, evaluating their capabilities in both chain-of-thought and tool-augmented reasoning on benchmarks in English and Chinese. Key observations are as follows: (1) When employing chain-of-thought reasoning, DeepSeekMath-RL 7B achieves accuracy scores of 88.2\% on GSM8K and 51.7\% on MATH. This not only outperforms every open-source model within the 7B to 70B parameter scale but also exceeds the performance of most proprietary systems. (2) Notably, DeepSeekMath-RL 7B’s training is confined exclusively to chain-of-thought instructional data derived from GSM8K and MATH, and its initial checkpoint is DeepSeekMath-Instruct 7B. Despite such limited training data, the model consistently surpasses DeepSeekMath-Instruct 7B across all assessed metrics, demonstrating the substantial gains conferred by reinforcement learning.

\section{Discussion}
In the following section, we present insights obtained from both pre-training and reinforcement learning experiments.

\subsection{Lessons Learnt in Pre-Training}

We begin by detailing our pre-training insights. Unless indicated otherwise, experimental configurations align with those provided in Section \ref{sec:quality-policy}. It should be emphasized that references to the DeepSeekMath Corpus in this discussion pertain to an 89B-token dataset, produced during the second round of data acquisition.

\subsubsection{Code Training Benefits Mathematical Reasoning}

There is a widely held but insufficiently substantiated belief that training with code improves models’ reasoning capabilities. We provide empirical evidence specific to mathematical tasks: integrating code in training enhances mathematical reasoning performance in scenarios both with and without tool assistance.

To examine the impact of code training on mathematical reasoning, we designed experiments involving two types of training regimes: two-stage training and one-stage training.

\textbf{Two-Stage Training}

\begin{itemize}[topsep=0pt]
    \item \textbf{Code Training for 400B Tokens $\rightarrow$ Math Training for 150B Tokens}: DeepSeek-LLM 1.3B is first pretrained with 400B code tokens, and then further trained with 150B tokens focused on mathematics.
    \item \textbf{General Training for 400B Tokens $\rightarrow$ Math Training for 150B Tokens}: For comparison, we conduct a control experiment where, in the initial phase, the model is pretrained using 400B general tokens sampled from a large, diverse corpus assembled by DeepSeek-AI, rather than code-specific data. This approach aims to evaluate whether code tokens offer unique advantages in developing mathematical reasoning compared to broadly sourced text.
\end{itemize}

\textbf{One-Stage Training}

\begin{itemize}[topsep=0pt]
    \item \textbf{Math Training for 150B Tokens}: The model, DeepSeek-LLM 1.3B, is trained exclusively on 150B mathematical tokens.
    \item \textbf{Training on a mixture of 400B Code Tokens and 150B Math Tokens}: Sequential math training after code pretraining can impair coding capabilities. Therefore, we examine whether joint, one-stage training on a blend of 400B code and 150B math tokens both enhances mathematical reasoning and lessens the detrimental impact of catastrophic forgetting typically seen in two-stage paradigms.
\end{itemize}

\paragraph{Results}
The performance across different training regimes is summarized in Table \ref{tab:code-math} and Table \ref{tab:code-math-general-eval}.

Incorporating code data during pretraining supports improved program-aided mathematical reasoning in both sequential and unified training protocols. According to the findings in Table \ref{tab:code-math}, sequentially training on code tokens dramatically elevates the model’s proficiency in tasks such as GSM8K and MATH when employing Python-based problem-solving, and the subsequent mathematical fine-tuning further boosts these capabilities. Notably, joint training that interleaves code and math tokens in a single stage not only addresses the catastrophic forgetting observed in two-stage schemes but also enhances both code-related tasks (Table \ref{tab:code-math-general-eval}) and math problems solvable with programmatic assistance (Table \ref{tab:code-math}).

Beyond program-aided tasks, code-focused training bolsters mathematical reasoning abilities even without tool usage. In two-stage settings, code pretraining provides moderate yet meaningful performance gains and primes the model for more effective math specialization, ultimately yielding the strongest results. In contrast, mixing code and math tokens in a one-stage process can diminish pure mathematical reasoning not involving code; this may be attributed to DeepSeek-LLM 1.3B’s limited size, restricting its ability to internalize both coding and math domains concurrently.

\subsubsection{ArXiv Papers Appear Unhelpful for Enhancing Mathematical Reasoning}

Despite being a prevalent element within mathematical pre-training datasets \citep{minerva,gpt-f,llemma,mathpile}, the direct contribution of arXiv papers to mathematical reasoning performance has not been thoroughly scrutinized. Somewhat unexpectedly, our empirical results indicate that arXiv papers offer limited benefits for mathematical reasoning skills. Our experiments leverage models across different parameter scales—specifically, DeepSeek-LLM 1.3B and DeepSeek-Coder-Base-v1.5 7B \citep{deepseek-coder}—using arXiv-based corpora subject to diverse preprocessing procedures:

\begin{itemize}[topsep=0pt]
    \item \textbf{MathPile} \citep{mathpile}: an 8.9B-token dataset extensively cleaned and filtered using heuristic criteria, predominantly composed (over 85\%) of scientific arXiv papers;
    \item \textbf{ArXiv-RedPajama} \citep{redpajama}: the complete collection of arXiv LaTeX sources with extraneous elements like preambles, comments, macros, and bibliographies excluded, totaling 28.0B tokens.
\end{itemize}

Our approach involves training DeepSeek-LLM 1.3B for 150B tokens and DeepSeek-Coder-Base-v1.5 7B for 40B tokens on each arXiv corpus separately. Results suggest that, when models are exposed solely to arXiv-derived data, there is little to no positive effect—sometimes even regression—in mathematical reasoning capabilities, as assessed by an array of benchmarks of varying difficulty levels in this study. These benchmarks include datasets focused on quantitative reasoning, such as GSM8K and MATH (see Table \ref{tab:arxiv-cot}), multiple-choice assessments like MMLU-STEM (Table \ref{tab:arxiv-cot}), and formal mathematics tasks such as miniF2F (Table \ref{tab:arxiv_atp}).

Nevertheless, these findings are subject to several caveats and should be interpreted with caution. We have not yet examined:

\begin{itemize}[topsep=0pt]
    \item The effect of arXiv-based data on specialized mathematical tasks beyond this paper’s scope, for example, the informalization of theorems, i.e., translating formal statements or proofs into informal discourse;
    \item The impact arising from the interplay between arXiv data and other dataset types;
    \item The possibility that benefits from arXiv content emerge only at larger model scales.
\end{itemize}

As such, more comprehensive investigations are warranted, which we defer to subsequent research.

\subsection{Insights of Reinforcement Learning}
\subsubsection{Toward a Unified Paradigm} Here, we develop a unified framework for the analysis of training paradigms—including SFT, RFT, DPO, PPO, and GRPO—and conduct experiments to investigate the underlying determinants of this framework. In broad terms, the gradient with respect to model parameters $\theta$ across these training algorithms can be expressed as:

\begin{equation}
    \nabla_{\theta}\mathcal{J}_{\textcolor{red}{\mathcal{A}}}(\theta) = \mathbb{E}[\underbrace{(q,o) \sim \textcolor{red}{\mathcal{D}}}_{Data \ Source}]\left( \frac{1}{|o|} \sum_{t=1}^{|o|}  \underbrace{GC_{{\mathcal{A}}}(q, o, t, \textcolor{red}{\pi_{{rf}}})}_{Gradient \ Coefficient}  \nabla_{\theta}\log \pi_{\theta}(o_t | q, o_{<t})\right).
\label{eq:objective}
\end{equation}

This formulation comprises three principal elements: 1) \textit{Data Source $\mathcal{D}$}, which specifies the dataset used during training; 2) \textit{Reward Function $\pi_{{rf}}$}, responsible for providing the reward signal guiding the learning process; 3) \textit{Algorithm $\mathcal{A}$}, which transforms both the training data and the associated reward signal into the gradient coefficient $GC$, thereby determining the extent to which each data point is penalized or reinforced. The following summarizes various prominent approaches within this overarching framework:

\begin{itemize}
    \item  \textbf{Supervised Fine-tuning (SFT)}: In SFT, the pretrained model is further adapted using human-curated SFT datasets.
    \item \textbf{Rejection Sampling Fine-tuning (RFT)}: RFT takes the SFT-trained model and continues fine-tuning on outputs, sampled in response to SFT prompts, that pass a filtering criterion based on answer accuracy.
    \item \textbf{Direct Preference Optimization (DPO)}: DPO builds on the SFT model by fine-tuning with additional outputs generated by the SFT model, leveraging a pairwise DPO loss to guide learning.
    \item \textbf{Online Rejection Sampling Fine-tuning (Online RFT)}: Unlike standard RFT, Online RFT initiates with the SFT model as the policy and updates it iteratively, using freshly sampled outputs from the current version of the policy in real time.
    \item \textbf{PPO/GRPO}: PPO/GRPO also employs the SFT model as the starting policy, but uses reinforcement learning with outputs dynamically sampled from the live policy to enhance performance.
\end{itemize}

The elements constituting these methods are outlined in Table \ref{tab:data_policy}. For a thorough derivation and more comprehensive analysis, please consult Appendix \ref{app:analysis-rl}.

\paragraph{Observation about Data Source} We categorize data sources into two main types: online sampling and offline sampling. Online sampling refers to collecting training data dynamically from the ongoing exploration trajectories of the current policy during training. In contrast, offline sampling relies on data obtained from sampling the initial SFT model before training commences. Both RFT and DPO employ the offline sampling paradigm, whereas Online RFT and GRPO adopt an online approach.

Referring to Figure \ref{fig:rl-analysis}, we observe that Online RFT demonstrates a pronounced improvement over RFT across both evaluation benchmarks. While their performance remains similar during the initial training phase, Online RFT quickly establishes a marked lead as training advances, highlighting the effectiveness of online learning. This phenomenon can be rationalized by noting that, at the outset, the actor closely resembles the SFT model and hence generates samples with minimal deviation. As training progresses, however, the actor produces samples with increasingly divergent characteristics, and continuous data collection from the evolving policy confers substantial benefits.

\paragraph{Observation about Gradient Coefficient} The transformation of the reward signal into a gradient coefficient is crucial for updating the model parameters. In our experimental framework, the reward function is classified as either `Rule'-based or `Model'-based. The `Rule' category involves assessing response quality by checking the factual accuracy of answers, while the `Model' category utilizes a reward model, trained on data labeled by rule-based assessment, to assign scores to responses. Equations \ref{eq:GC-RFT} and \ref{eq:GC-GRPO} reveal a fundamental distinction: GRPO adapts the gradient coefficient in accordance with the reward magnitude output by the reward model, thus enabling differential scaling of reinforcement and penalization based on response quality. Conversely, Online RFT does not implement this nuance; it fails to penalize erroneous responses and reinforces all correct answers uniformly, regardless of their individual merit.

As illustrated in Figure \ref{fig:rl-analysis}, GRPO outperforms online RFT, underscoring the effectiveness of dynamically adjusting positive and negative gradient coefficients. Moreover, GRPO+PS demonstrates enhanced results relative to GRPO+OS, which suggests the advantage of adopting step-aware, fine-grained gradient coefficients. Our investigation extends to iterative RL, where two rounds of iteration were performed. As depicted in Figure \ref{fig:iter-rl}, iterative RL produces notable performance gains, particularly during the initial iteration.

\subsubsection{Why RL Works?}  
In this study, reinforcement learning is applied on a subset of the instruction tuning dataset and yields substantial improvements over the instruction-tuned baseline. To further probe the mechanisms underlying the effectiveness of reinforcement learning, we compare the Pass@K and Maj@K accuracies of both the Instruct and RL models across two benchmarks. As seen in Figure \ref{fig:maj-pass}, RL primarily boosts Maj@K, with little to no effect on Pass@K. These observations imply that the benefits from RL are mainly due to enhancing the resilience of the output distribution—in particular, \textbf{the gains appear to arise from increasing the likelihood that at least one correct response is found among the TopK outputs, as opposed to improving the core modeling capacity.} This aligns with findings in \citep{wang2023making}, who characterized a \textbf{misalignment problem} for reasoning tasks in SFT models and demonstrated that applying preference alignment strategies could yield better reasoning performance \citep{yuan2023rrhf,song2023preference,wang2023making}.

\subsubsection{How to Achieve More Effective RL?}
\label{sec:effec_rl}
Our experiments confirm the effectiveness of RL in mathematical reasoning contexts, and we put forth a unifying framework that conceptualizes a variety of prevalent training methods as forms of direct or simplified RL. Within this framework, the three essential elements are: Data Source, Algorithm, and Reward Function, as delineated in Equation \ref{eq:objective}. Several promising avenues for further advancement across these components are discussed below.

\paragraph{Data Source}  
The choice of data source is foundational for any training paradigm. For RL, this specifically refers to unlabeled questions alongside outputs generated by sampling from the policy model. This study limits the data to questions used during instruction tuning and employs simple nucleus sampling for response generation—a limitation we suspect constrains the improvements primarily to Maj@K. Looking forward, we intend to extend our RL methods to out-of-distribution questions, and pair them with \textbf{more sophisticated sampling (decoding) techniques}, such as those leveraging tree-search algorithms \citep{yao2023tree}. Additionally, emerging \textbf{efficient inference methods} \citep{xia-etal-2023-speculative,leviathan2023fast,kwon2023efficient,xia2024unlocking}, which critically influence the efficiency of policy model exploration, are poised to become increasingly impactful.

\paragraph{Algorithms} Algorithms use both data and reward signals to calculate gradients, which are then applied to adjust the parameters of the model. As indicated by Equation \ref{eq:objective}, current methodologies almost invariably place full \textbf{TRUST} in the reward signal, modifying the conditional probability of individual tokens accordingly. Nevertheless, such confidence is not always justified, particularly when faced with highly intricate tasks, where the fidelity of the reward signal may be questionable. For instance, even with the PRM800K datasets \citep{lightman2023let}, which have benefited from meticulous annotation by skilled annotators, about 20\% of the labels remain incorrect\footnote{\url{https://github.com/openai/prm800k/issues/12\#issuecomment-1728491852}}. Therefore, it becomes essential to investigate reinforcement learning approaches that are resilient in the face of noisy or unreliable reward feedback. We contend that alignment frameworks following the \textbf{WEAK-TO-STRONG} paradigm \citep{burns2023weak} have the potential to revolutionize the core of learning algorithms.

\paragraph{Reward Function} The reward function constitutes the primary signal for model training. In reinforcement learning contexts, the reward function often takes the form of a neural reward model. We identify three critical avenues for advancing reward model design: 1) \textbf{Improving the generalization capacity of reward models.} Reward models must possess robust generalization to manage unseen queries and complex decoding outputs; failing this, reinforcement learning risks merely anchoring LLM distributions without enhancing their intrinsic performance; 2) \textbf{Capturing and expressing uncertainty in the reward model.} Accurately representing uncertainty may provide the necessary link between underpowered reward models and weak-to-strong alignment algorithms; 3) \textbf{Developing effective high-quality process reward models} that deliver granular, step-by-step feedback during reasoning processes \citep{lightman2023let,wang2023math}.

\section{Conclusion, Limitation, and Future Work} In this work, we introduce DeepSeekMath, which sets a new standard among open-source models on the competitive MATH benchmark, achieving results that closely rival those of proprietary systems. DeepSeekMath~builds upon DeepSeek-Coder-v1.5 7B and is subjected to an extended training regimen comprising 500B tokens, notably incorporating 120B mathematical tokens curated from Common Crawl as a substantial portion of its dataset. Our comprehensive ablation analyses highlight the considerable value of web content in procuring quality mathematical materials, whereas arXiv sources appear to contribute less than initially anticipated. We propose Group Relative Policy Optimization (GRPO), a novel adaptation of Proximal Policy Optimization (PPO), which demonstrably enhances mathematical reasoning while utilizing less memory. Empirical findings confirm that GRPO remains beneficial even after DeepSeekMath-Instruct 7B attains strong benchmark performance. Additionally, we offer a unified framework for interpreting various methods and outline several avenues for advancing reinforcement learning methodologies.

Despite the strong results achieved by DeepSeekMath~on benchmarks involving quantitative reasoning, it remains less proficient than proprietary models in areas such as geometry and theorem proving. During preliminary testing, the model exhibited challenges with tasks involving triangles and ellipses, suggesting possible data distribution imbalances during pre-training and fine-tuning stages. Furthermore, given current model size limitations, DeepSeekMath~lags behind GPT-4 in few-shot performance; GPT-4 shows clear improvements when provided with few-shot examples, whereas DeepSeekMath~performs comparably in both zero-shot and few-shot evaluations. Looking ahead, we plan to refine our data selection strategies to develop an even higher quality pre-training dataset. Moreover, we aim to investigate the proposed potential directions (Section \ref{sec:effec_rl}) to facilitate more effective reinforcement learning strategies for large language models.

\end{document}